% =====================================================================
%  Chapter 35 — The Taxable Investor's Record
% =====================================================================

\chapter{The Taxable Investor's Record}
\label{ch:taxable-record}

The system's tax-exempt record is the cleanest thing the book measures. The
taxable record is the one the reader is most likely to have, and it is
worse. This chapter reports the difference, explains where it comes from,
and draws the conclusion that a reader who has a taxable account\idx{Taxable account} must act
on.

\section{The two accounts}

The system runs two portfolios in parallel from the same month, with the
same weights and the same trades. The only difference is that one pays
tax.

\begin{center}
\begin{tabular}{@{}P{1.70in}rrrr@{}}
\toprule
 & Compound growth & Sharpe & Deepest drawdown & Worst 12 months \\
\midrule
Tax-exempt account\idx{Tax-exempt account} & \numFullCAGR{} & \numFullSharpe{} & \numFullDD{} & \numFullWorstYear{} \\
Taxable account    & \numTaxCAGR{} & \numTaxSharpe{} & \numTaxDD{} & \numTaxWorstYear{} \\
\bottomrule
\end{tabular}
\end{center}

The cost is \numTaxPenalty{} a year of compound growth. Over the
\numRecordMonths{} months of the record, the taxable account pays
\numTaxTotal{} dollars of tax on a starting capital of \numInitialCapital{}
dollars --- \numTaxFraction{} of the initial capital, paid away.

The two accounts, and the wedge between them, are shown in \figref{fig:tax}.

\figref{fig:tax} makes a point that is easy to miss: the wedge is
not constant. It widens, because the tax is charged on a portfolio whose
value is growing, and because the system's turnover keeps realising gains
into a rising base.

\section{Where the money goes}

Three components, and it is worth separating them because they have
different remedies.

\textbf{The preferential rate is good news.} Long-term capital gains are
taxed at fifteen per cent for most taxpayers, against a marginal ordinary
rate that is typically thirty-two or thirty-seven per cent. A portfolio
that trades a great deal and holds long-term gains is, on this account,
better off than a buy-and-hold investor in a taxable account, because the
buy-and-hold investor eventually realises a large long-term gain\idx{Long-term gain} and pays
the tax on it. The book's system realises gains every year, at the
preferential rate, and the wedge it creates is smaller than the wedge a
liquidation event would create.

\textbf{Short-term gains are the expensive part.} A gain realised within a
year is taxed as ordinary income, at the full marginal rate plus the net
investment income tax. The book's no-trade band\idx{No-trade band} and the system's low
turnover --- \numTurnover{} a month --- keep most of the realised gains in
the long-term bucket. This is the single largest design decision in the
tax layer: a system that rebalanced more aggressively would convert
long-term rates into short-term ones and would be materially worse after
tax.

\textbf{The wash-sale rule is the tax layer's own constraint.} Because
the rule applies in both directions, the system cannot harvest a loss and
immediately buy the same fund back. Over the record, this occasionally
prevents the system from doing something it would otherwise have done ---
holding a losing position longer than it would like --- and the effect is
a small drag on the taxable account relative to a naive backtest that
ignored the rule.

\section{The mechanical result worth stating}

There is a comparison worth making, because it is the argument for
automated, patient, long-horizon management in a taxable account.

The taxable system earns \numTaxCAGR{} with a drawdown of \numTaxDD{},
while the benchmark earns \numSpyCAGR{} with a drawdown of \numSpyDD{} in
an account with no tax at all. But a real investor in the benchmark does
not experience \numSpyCAGR{}: at the end, they own the whole position and
must realise it, and the tax on that realisation is charged on the
entire accumulated gain.

The book does not model that liquidation --- it would be a single
transaction and the rates are known --- but it is worth naming the
asymmetry. A system that realises gains gradually, at the preferential
rate, every year, converts a deferred tax liability into an annually
paid, annually-priced one. The deferral is a benefit to the holder of the
index that the system does not have to pay for.

\begin{takeaway}[The taxable conclusion]
The tax costs \numTaxPenalty{} a year of compound growth, and it is paid
on a portfolio that keeps growing. The system's low turnover is what
keeps most of the realised gains in the long-term bucket, and that is a
design decision rather than an accident: the no-trade band is a tax
instrument as much as a cost instrument.
\end{takeaway}

\section{What this means for a reader}

If the reader has a taxable account and a long horizon, the conclusion is
that the tax layer is not a detail to be deferred and it is not a reason
to abandon the approach. It is a reason to make the system's turnover low,
to hold the tax-exempt account for the assets that are not being traded,
and to accept that the tax-exempt record is the honest measure of the
strategy and the taxable record is the honest measure of the investor's
outcome.

If the reader has a taxable account and a short horizon, the conclusion
is the opposite: the tax friction compounds with the behavioural friction
of Chapter 8, and a system that rebalances is likely to be the wrong
instrument.

\begin{keytakeaways}
\item The taxable account pays \numTaxTotal{} dollars over the record,
\numTaxFraction{} of the initial capital, and gives up \numTaxPenalty{} a
year of compound growth.
\item The tax is charged on a growing portfolio, so the wedge widens.
\item Low turnover keeps most realised gains in the long-term bucket; the
no-trade band is a tax instrument as much as a cost instrument.
\item A taxable reader with a short horizon should not run a rebalancing
system at all.
\end{keytakeaways}
